% Options for packages loaded elsewhere
\PassOptionsToPackage{unicode}{hyperref}
\PassOptionsToPackage{hyphens}{url}
\PassOptionsToPackage{dvipsnames,svgnames,x11names}{xcolor}
\documentclass[
  11pt,
]{article}
\usepackage{xcolor}
\usepackage[margin=2.5cm]{geometry}
\usepackage{amsmath,amssymb}
\setcounter{secnumdepth}{-1}
\usepackage{iftex}
\ifPDFTeX
  \usepackage[T1]{fontenc}
  \usepackage[utf8]{inputenc}
  \usepackage{textcomp} % provide euro and other symbols
\else % if luatex or xetex
  \usepackage{unicode-math} % this also loads fontspec
  \defaultfontfeatures{Scale=MatchLowercase}
  \defaultfontfeatures[\rmfamily]{Ligatures=TeX,Scale=1}
\fi
\usepackage{lmodern}
\ifPDFTeX\else
  % xetex/luatex font selection
\fi
% Use upquote if available, for straight quotes in verbatim environments
\IfFileExists{upquote.sty}{\usepackage{upquote}}{}
\IfFileExists{microtype.sty}{% use microtype if available
  \usepackage[]{microtype}
  \UseMicrotypeSet[protrusion]{basicmath} % disable protrusion for tt fonts
}{}
\makeatletter
\@ifundefined{KOMAClassName}{% if non-KOMA class
  \IfFileExists{parskip.sty}{%
    \usepackage{parskip}
  }{% else
    \setlength{\parindent}{0pt}
    \setlength{\parskip}{6pt plus 2pt minus 1pt}}
}{% if KOMA class
  \KOMAoptions{parskip=half}}
\makeatother
\usepackage{longtable,booktabs,array}
\newcounter{none} % for unnumbered tables
\usepackage{calc} % for calculating minipage widths
% Correct order of tables after \paragraph or \subparagraph
\usepackage{etoolbox}
\makeatletter
\patchcmd\longtable{\par}{\if@noskipsec\mbox{}\fi\par}{}{}
\makeatother
% Allow footnotes in longtable head/foot
\IfFileExists{footnotehyper.sty}{\usepackage{footnotehyper}}{\usepackage{footnote}}
\makesavenoteenv{longtable}
\setlength{\emergencystretch}{3em} % prevent overfull lines
\providecommand{\tightlist}{%
  \setlength{\itemsep}{0pt}\setlength{\parskip}{0pt}}
\usepackage{bookmark}
\IfFileExists{xurl.sty}{\usepackage{xurl}}{} % add URL line breaks if available
\urlstyle{same}
\hypersetup{
  pdftitle={ILCSS --- Iterative Longest Common SubStrings Similarity: A String Similarity Measure for Multi-Part Personal Name Matching},
  pdfauthor={Raoul Sotero Calao},
  colorlinks=true,
  linkcolor={Maroon},
  filecolor={Maroon},
  citecolor={Blue},
  urlcolor={Blue},
  pdfcreator={LaTeX via pandoc}}

\title{ILCSS --- Iterative Longest Common SubStrings Similarity: A
String Similarity Measure for Multi-Part Personal Name Matching}
\author{Raoul Sotero Calao\thanks{\footnotesize\begin{tabular}[t]{@{}l@{}}
  Chief Commissioner, Italian National Police - Central Anticrime Directorate\\
  (\texttt{raoulsotero.calao@poliziadistato.it}; \texttt{rscalao@gmail.com})
\end{tabular}}}
\date{2026\\[0.5em]\small Preprint: \href{https://doi.org/10.5281/zenodo.19462182}{doi:10.5281/zenodo.19462182}}

\begin{document}
\maketitle
\begin{abstract}
We present ILCSS (Iterative Longest Common SubStrings Similarity), a
string similarity measure specifically designed for approximate matching
of multi-part personal names from heterogeneous sources. The problem
arises when the same individual's name appears with its components in
different orders or with a different number of parts across
administrative databases, criminal intelligence systems, or
bibliographic records --- a case that defeats character-level measures
such as Levenshtein distance and Jaro-Winkler. ILCSS iteratively
extracts the longest common substrings shared by two strings, consuming
matched blocks at each step and applying a positional penalty that
attenuates with block length, thereby rewarding component-level
agreement while retaining sensitivity to ordering conventions. We prove
that ILCSS is a similarity coefficient --- bounded, reflexive, and
conditionally symmetric --- but not a metric distance. A synthetic
benchmark of 295 name pairs drawn from five naming traditions (Arabic,
Hispanic, Portuguese, Italian, French/German), covering all
permutations, component-drop, and insertion variants, shows that ILCSS
achieves a same-person/different-person score gap of 0.729, compared to
0.262--0.361 for competing measures, with F1 = 0.998 at an optimal
threshold of 0.40. Evaluation on two real-world record-linkage datasets
(DBLP-ACM and FEBRL) confirms the advantage, with F1 = 0.994 and 0.989
respectively. Reference implementations in Perl, Crystal, and Python are
provided.
\end{abstract}

\section{1. Problem Statement}\label{problem-statement}

Matching personal names across heterogeneous databases is a central
challenge in record linkage and entity resolution {[}18, 19{]}. The
difficulty goes beyond simple typographic variation: names from
different sources --- civil registries, immigration records,
administrative databases, transliterated documents --- frequently differ
in the \textbf{order} and \textbf{number} of their constituent parts,
while the parts themselves remain identical or near-identical.

Consider the following pairs, each referring to the same individual:

{\def\LTcaptype{none} % do not increment counter
\begin{longtable}[]{@{}
  >{\raggedright\arraybackslash}p{(\linewidth - 4\tabcolsep) * \real{0.3333}}
  >{\raggedright\arraybackslash}p{(\linewidth - 4\tabcolsep) * \real{0.3333}}
  >{\raggedright\arraybackslash}p{(\linewidth - 4\tabcolsep) * \real{0.3333}}@{}}
\toprule\noalign{}
\begin{minipage}[b]{\linewidth}\raggedright
Source A
\end{minipage} & \begin{minipage}[b]{\linewidth}\raggedright
Source B
\end{minipage} & \begin{minipage}[b]{\linewidth}\raggedright
Variation type
\end{minipage} \\
\midrule\noalign{}
\endhead
\bottomrule\noalign{}
\endlastfoot
\texttt{mario\ rossi} & \texttt{rossi\ mario} & Two-part
transposition \\
\texttt{mario\ rossi} & \texttt{mario\ giuseppe\ rossi} & Middle name
insertion \\
\texttt{abdullah\ abdel\ aziz} & \texttt{aziz\ abdullah\ abdel} &
Three-part permutation \\
\texttt{abdullah\ abdel\ aziz} & \texttt{abdel\ aziz\ abdullah} &
Three-part permutation \\
\texttt{maria\ rosa\ garcia\ lopez} &
\texttt{garcia\ lopez\ maria\ rosa} & Block transposition \\
\end{longtable}
}

This problem is particularly acute for Arabic names, which commonly
consist of a chain of three or more components --- a given name followed
by the father's name, grandfather's name, and optionally a family name
--- recorded in different orders across different administrative
traditions. It also arises systematically with Hispanic double-surname
naming conventions, Portuguese and Brazilian multi-part surnames, and
historical records where name order was not standardized.

A practically important domain where this challenge is especially severe
is \textbf{law enforcement and criminal intelligence analysis}.
Investigations targeting organized crime networks and domestic or
international terrorism routinely require cross-referencing personal
name data across heterogeneous sources --- criminal records, immigration
and border control systems, intelligence reports, court documents,
watchlists, and international law-enforcement information exchanges such
as Interpol notices or Europol databases --- each of which may follow
different national conventions for name ordering and may record a
different subset of an individual's name components. A subject who
appears as \texttt{"ali\ hassan\ abdel\ rahman"} in one system may be
recorded as \texttt{"abdel\ rahman\ ali"} or simply
\texttt{"hassan\ ali"} in another, and standard character-level
similarity measures will fail to identify these as the same individual.
The consequences of a missed match --- or of a spurious one --- in this
context are operationally significant, making reliable approximate name
matching a genuine requirement rather than a convenience.

Existing string similarity measures handle these cases poorly.
\textbf{Levenshtein edit distance} {[}3{]} treats each character
substitution, insertion, and deletion as a unit cost: transposing two
long name parts requires rewriting almost the entire string, producing a
large distance even when all parts are identical. \textbf{Jaro-Winkler}
{[}15, 16{]} explicitly rewards agreement at the beginning of the string
and applies a positional window for character matching; for
\texttt{"abdullah\ abdel\ aziz"} vs \texttt{"aziz\ abdullah\ abdel"},
neither the prefix nor the character-position assumptions hold, and the
score is correspondingly low. \textbf{Ratcliff/Obershelp} {[}2{]} finds
the longest common block recursively and computes similarity as twice
the matched characters over total characters; it correctly identifies
the common blocks but has no mechanism to distinguish a same-person pair
--- where all parts are present, merely reordered --- from a
different-person pair where only some parts coincide.

\textbf{Token-based methods} (Jaccard on word tokens, sorted-token
similarity) address the ordering problem by treating the name as a bag
of words, but discard all information about the relative positioning of
parts, which carries genuine signal when sources are expected to use
similar conventions.

\begin{center}\rule{0.5\linewidth}{0.5pt}\end{center}

\section{2. The ILCSS Approach}\label{the-ilcss-approach}

ILCSS (Iterative Longest Common SubStrings) measures the similarity
between two strings in the range {[}0, 1{]}. It treats the full name ---
including all parts and spaces --- as a single character string. Rather
than requiring a single global alignment, it \textbf{iterates}: at each
step it identifies all longest common substrings shared by the two
strings, ``consumes'' them by replacing them with placeholders (so they
are not re-matched), and adds their contribution to the score. The
process repeats until no further block of length $\geq$ 3 remains.

The key design element is the \textbf{positional penalty}: each matched
block contributes to the score in proportion to its length, but the
contribution is attenuated if the block appears at very different
positions in the two strings. The attenuation itself decreases for
longer blocks --- a long block matching across a large positional gap is
still strong evidence of identity. This balances two competing signals:
same-part content (matched block length) and same-ordering convention
(positional agreement).

For the pair \texttt{"abdullah\ abdel\ aziz"} /
\texttt{"aziz\ abdullah\ abdel"}:

\begin{enumerate}
\def\labelenumi{\arabic{enumi}.}
\tightlist
\item
  First iteration finds \texttt{"abdullah"} (length 8) in both strings,
  at positions 0 and 5 respectively. The positional distance is 5 out of
  L = 19; the penalty is moderate and attenuated by the block's relative
  length (8/19 $\approx$ 0.42). Block score $\approx$ 7.1.
\item
  After consuming \texttt{"abdullah"}, finds \texttt{"abdel"} (length 5)
  at positions 9 and 12 in the modified strings. Similar attenuation.
  Block score $\approx$ 4.3.
\item
  After consuming \texttt{"abdel"}, finds \texttt{"aziz"} (length 4) at
  positions 15 and 0. Larger positional gap, but again attenuated. Block
  score $\approx$ 2.9.
\item
  No further blocks of length $\geq$ 3. Total score $\approx$ 14.3. Similarity $\approx$ 14.3
  / 19 $\approx$ \textbf{0.75}.
\end{enumerate}

The same pair scored by Levenshtein (normalized) gives approximately
0.37; Jaro-Winkler gives approximately 0.55. ILCSS correctly identifies
the pair as highly similar.

\begin{center}\rule{0.5\linewidth}{0.5pt}\end{center}

\section{3. Pseudocode}\label{pseudocode}

\begin{verbatim}
function ilcss_sim(s1, s2):
    s1 := remove_spaces(s1)
    s2 := remove_spaces(s2)
    L := max(length(s1), length(s2))      # reference length for normalization
    total_score := 0

    repeat:
        blocks := longest_common_substrings(s1, s2)
                  # returns all common substrings of maximum length
        min_len := length(blocks[0])      # all blocks have the same length

        for each block in blocks:
            b := length(block)
            if b < 3: skip to next       # blocks shorter than a trigram are ignored

            p1 := position of block in s1
            p2 := position of block in s2

            # "consume" the block so it is not matched in subsequent iterations
            s1 := replace(s1, p1, b, "{" * b)
            s2 := replace(s2, p2, b, "}" * b)

            # penalty based on positional distance between the two blocks
            dist        := |p2 - p1|
            r_dist      := dist / L                # relative distance
            r_len       := b / L                   # relative block length
            r_penalty   := r_dist * (1 - r_len)   # longer blocks reduce the penalty
            total_score := total_score + b * (1 - r_penalty)

    until min_len <= 2

    return total_score / L               # similarity normalized in [0, 1]
\end{verbatim}

\begin{verbatim}
function longest_common_substrings(s1, s2):
    # space-optimized DP [3, 5]: only two columns kept in memory
    prev[0..m] := 0,  curr[0..m] := 0    # m = length of s2
    best := 0
    results := []

    for i from 1 to length(s1):
        for j from 1 to length(s2):
            if s1[i-1] == s2[j-1]:            # 0-indexed character access
                curr[j] := prev[j-1] + 1
                if curr[j] > best:
                    best := curr[j]
                    results := [ s1[i-best .. i] ]       # new list with this one only
                else if curr[j] == best:
                    results.append( s1[i-best .. i] )   # tie: add to the list
            else:
                curr[j] := 0
        swap prev and curr
        reset curr to 0

    return results
\end{verbatim}

\begin{center}\rule{0.5\linewidth}{0.5pt}\end{center}

\section{4. Relationship to prior
work}\label{relationship-to-prior-work}

\subsection{4.1 Record linkage and personal name
matching}\label{record-linkage-and-personal-name-matching}

The probabilistic framework for record linkage was established by
Fellegi \& Sunter {[}18{]}, who formalized the problem of determining
whether two records from different sources refer to the same entity.
Personal name matching is among the most studied subtasks within this
framework {[}19{]}. Surveys of name-matching techniques --- notably
Christen {[}20{]} and Lait \& Randell {[}21{]} --- consistently identify
the core difficulty as the mismatch between the diversity of real-world
name representations and the assumptions of character-level distance
measures.

The particular challenge of \textbf{multi-part names with variable
component order} has received growing attention in several specific
domains. For Spanish names, Ruiz-P\'erez et al.~{[}22{]} documented that
between 48\% and 69\% of Spanish authors appear under multiple name
forms across major bibliographic databases, largely because of
inconsistent handling of the double-surname convention; no dedicated
matching algorithm for this problem is available in the literature. For
genealogical records, Wilson {[}23{]} explicitly addresses
first-name/last-name transposition and proposes combined-field indexing
as a mitigation, without resolving the general case of arbitrary
component permutations.

Arabic personal names pose the most severe instance of this problem.
Arabic naming tradition commonly chains three or more components ---
given name, father's name, grandfather's name, and optionally a tribal
or family name --- and different administrative systems record these
chains in different orders. El-Shishtawy {[}24{]} addresses Arabic
name-based database linkage directly, proposing a rule-based system for
name normalization prior to matching. Freeman, Condon \& Ackerman
{[}25{]} extend edit distance with Arabic-English character equivalency
classes to handle transliteration variants, but neither work addresses
the component-reordering problem algorithmically. Hassanat \& Altarawneh
{[}26{]}, working from a corpus of 9.9 million Jordanian citizen names,
systematically catalogue the orthographic and phonological spelling
variants that arise in Arabic name databases --- character
substitutions, Hamzah deletion, definite-article attachment,
feminization suffixes, and encoding-driven substitutions from legacy
7-bit systems --- and propose a rule-and-dictionary approach for
normalisation. Their work addresses within-token variation exclusively;
the across-token component-reordering problem (the same person's name
components appearing in different orders across records) is outside
their scope and is not addressed by any of the standard string
similarity measures surveyed in the literature.

\subsection{4.2 String similarity
measures}\label{string-similarity-measures}

The Jaro comparator {[}15{]} and its Winkler extension {[}16{]},
developed for the U.S. Census name-matching problem, explicitly reward
prefix agreement and apply a character-matching window proportional to
string length. These design choices are well-suited for typographic
variants of single-word names but work against multi-part name matching:
for \texttt{"abdullah\ abdel\ aziz"} vs
\texttt{"aziz\ abdullah\ abdel"}, prefix agreement is zero and the
positional window misaligns most character matches.

Token-based approaches --- Jaccard coefficient on word tokens,
sorted-token Jaro-Winkler (SoftTFIDF, as evaluated in Cohen, Ravikumar
\& Fienberg {[}17{]}) --- handle component reordering by treating the
name as a bag of words. This correctly addresses full permutations but
discards all positional information, so a name with all correct parts in
the wrong order scores identically to one where a part has been randomly
substituted. For a matching task where both same-person and
similar-but-different-person pairs are present, this loss of positional
signal is a liability.

The idea of iteratively extracting the longest common substring was
described as early as 1967 by Alberga {[}1{]}, who attributes it to
unpublished work by Baskin and Selfridge at IBM Research --- cited only
as private communications and never published independently. The closest
published algorithm sharing this iterative structure is the
Ratcliff/Obershelp method {[}2{]} (1988), which recursively finds the
longest common substring and then recurses on the flanking portions,
computing similarity as twice the number of matched characters divided
by total characters. ILCSS differs from Ratcliff/Obershelp in two key
respects: it applies a \textbf{positional penalty} (blocks found at very
different positions in the two strings contribute less to the score, but
less so for longer blocks), and it processes all blocks of equal maximum
length at each iteration rather than a single one.

The underlying dynamic programming recurrence for the longest common
substring is a standard result in string algorithmics {[}3, 4{]}. The
space optimization that keeps only two columns of the DP table in memory
--- reducing space from O(nm) to O(m) --- follows the approach described
by Hirschberg {[}5{]} for the related longest common subsequence
problem. For a broad survey of string similarity and approximate
matching techniques, see Navarro {[}6{]}.

\begin{center}\rule{0.5\linewidth}{0.5pt}\end{center}

\section{5. Key properties}\label{key-properties}

\textbf{Minimum block length.} Blocks shorter than 3 characters are
discarded. The minimum length of 3 reflects the well-established role of
trigrams as the smallest orthographic unit with sufficient
discriminative power for string matching. Blocks of length 1 or 2 match
too frequently by chance: over a 26-character Latin alphabet, there are
only 26$^2$ = 676 distinct bigrams, many of which (e.g., \emph{th},
\emph{he}, \emph{in}) occur with very high frequency in natural-language
text, producing a high rate of spurious matches. Trigrams yield 26$^3$ =
17,576 distinct types with a considerably flatter distribution,
providing meaningful selectivity without the storage and computational
overhead of longer n-grams {[}7{]}. This choice is consistent with the
information-theoretic argument underlying Ukkonen's q-gram filter
{[}8{]}: smaller values of q produce weaker filters, allowing more
candidate pairs to pass through. Empirically, Willett {[}9{]} directly
compared digram and trigram encoding on the Cranfield test collection
and found that \emph{``trigram encoding of words performs noticeably
better than the use of digrams''}; Angell, Freund and Willett {[}10{]}
subsequently adopted trigrams as the operative unit in their canonical
trigram-based spelling correction system, achieving over 90\% correct
identification against a 64,636-word dictionary.

\textbf{Iteration.} At each cycle the matched blocks are masked, so the
next iteration searches for matches in the remaining material, capturing
all shared name parts regardless of their number and order.

\textbf{Positional penalty.} A common block is worth less if it appears
at very different positions in the two strings. The formula
\texttt{r\_dist\ $\times$\ (1\ -\ r\_len)} attenuates the penalty for longer
blocks: a long name part is strong evidence of identity even if
positionally displaced. This means that fully reordered names (all parts
present, different order) score high, while partially overlapping names
(only some parts in common) score proportionally lower.

\textbf{Normalization.} Dividing by \texttt{L} (the length of the longer
string) guarantees a result in {[}0, 1{]}, where 1 means identical
strings (ignoring spaces) and 0 means no common substring of length $\geq$ 3.
Using the longer string as the reference penalizes cases where one name
is a proper subset of the other (e.g., \texttt{"mario\ rossi"} vs
\texttt{"mario\ giuseppe\ rossi"}).

\textbf{Complexity.} Each iteration runs the DP in O(n$\cdot$m); the number of
iterations is at most O(min(n,m)/3), so the overall complexity is
O(n$\cdot$m$\cdot$min(n,m)).

\begin{center}\rule{0.5\linewidth}{0.5pt}\end{center}

\section{6. Formal analysis of metric
properties}\label{formal-analysis-of-metric-properties}

ILCSS computes a \textbf{similarity measure} --- a function mapping
pairs of strings to {[}0, 1{]} --- not a distance. Following the
axiomatic framework of Santini \& Jain {[}11{]} and the survey by Lesot,
Rifqi \& Benhadda {[}12{]}, we examine each relevant property in turn.

\textbf{Boundedness and non-negativity.} Since
\texttt{r\_dist\ $\in$\ {[}0,1{]}} and \texttt{r\_len\ $\in$\ {[}0,1{]}}, we
have \texttt{r\_penalty\ =\ r\_dist\ $\times$\ (1\ --\ r\_len)\ $\in$\ {[}0,1{]}},
so \texttt{block\_score\ =\ b\ $\times$\ (1\ --\ r\_penalty)\ $\geq$\ 0}. The sum of
all matched block lengths cannot exceed
\texttt{min(\textbar{}s1\textbar{},\ \textbar{}s2\textbar{})\ $\leq$\ L},
hence \texttt{total\_score\ /\ L\ $\leq$\ 1}. Therefore \textbf{0 $\leq$ sim(s1,
s2) $\leq$ 1}. $\blacksquare$

\textbf{Reflexivity.} If \texttt{s1\ =\ s2} (after whitespace removal),
the first call to \texttt{longest\_common\_substrings} returns the whole
string as the unique block, with \texttt{dist\ =\ 0},
\texttt{r\_penalty\ =\ 0}, and
\texttt{block\_score\ =\ \textbar{}s1\textbar{}}. Thus
\texttt{total\_score\ /\ L\ =\ 1}. Conversely, \texttt{sim\ =\ 1}
implies matched length equals \texttt{L}, which --- given the no-overlap
constraint --- requires \texttt{s1\ =\ s2}. Hence \textbf{sim(s1, s2) =
1 $\Longleftrightarrow$ s1 = s2} (modulo spaces). $\blacksquare$

\textbf{Symmetry --- conditional.} The positional penalty
\texttt{\textbar{}p2\ --\ p1\textbar{}\ /\ L} and the normalization by
\texttt{L\ =\ max(\textbar{}s1\textbar{},\ \textbar{}s2\textbar{})} are
both symmetric under exchange of s1 and s2. However, when
\texttt{longest\_common\_substrings} returns multiple blocks of equal
maximum length, the tie-breaking order depends on the loop structure
(outer loop over s1), and subsequent \texttt{index} calls select the
first occurrence in the modified string. Swapping s1 and s2 may
therefore produce a different sequence of consumed blocks across
iterations and hence a different score. \textbf{Symmetry holds when the
LCS is unique at each iteration; it is not guaranteed in the presence of
ties.}

\textbf{Triangle inequality --- does not hold in general.} Converting
ILCSS to a distance via \texttt{d(s1,\ s2)\ =\ 1\ --\ sim(s1,\ s2)}, the
triangle inequality \texttt{d(s1,\ s3)\ $\leq$\ d(s1,\ s2)\ +\ d(s2,\ s3)} is
not satisfied in general. This follows from the result established by
Kondrak {[}13{]} that normalized LCS-based similarity measures lose the
triangle inequality when converted to distance: normalization by string
length introduces a non-linearity that invalidates the inequality. The
same phenomenon is documented for normalized edit distance by Li \& Liu
{[}14{]}, who propose a corrected normalization formula that restores
the metric property; an analogous correction for ILCSS is left as future
work.

\textbf{Positional sensitivity as a deliberate design choice.} The
positional penalty places ILCSS in the family of \emph{positionally
sensitive} string comparators, alongside the Jaro comparator {[}15{]}
and its Jaro-Winkler extension {[}16{]}, both of which are also
non-metrics. The underlying intuition --- that common substrings aligned
at similar positions provide stronger evidence of similarity than
displaced ones --- is empirically validated for name-matching and entity
resolution tasks in Cohen, Ravikumar \& Fienberg {[}17{]}.

\textbf{Summary.}

{\def\LTcaptype{none} % do not increment counter
\begin{longtable}[]{@{}
  >{\raggedright\arraybackslash}p{(\linewidth - 4\tabcolsep) * \real{0.3333}}
  >{\raggedright\arraybackslash}p{(\linewidth - 4\tabcolsep) * \real{0.3333}}
  >{\raggedright\arraybackslash}p{(\linewidth - 4\tabcolsep) * \real{0.3333}}@{}}
\toprule\noalign{}
\begin{minipage}[b]{\linewidth}\raggedright
Property
\end{minipage} & \begin{minipage}[b]{\linewidth}\raggedright
Holds?
\end{minipage} & \begin{minipage}[b]{\linewidth}\raggedright
Notes
\end{minipage} \\
\midrule\noalign{}
\endhead
\bottomrule\noalign{}
\endlastfoot
0 $\leq$ sim $\leq$ 1 & \checkmark{} & Proven above \\
sim(s, s) = 1 & \checkmark{} & Proven above \\
sim = 1 $\Longleftrightarrow$ s1 = s2 & \checkmark{} & Proven above \\
Symmetry & Conditional & Holds unless tie-breaking produces different
block sequences \\
Triangle inequality (1 -- sim) & $\times$ & Follows from Kondrak {[}13{]};
normalization destroys it \\
\end{longtable}
}

ILCSS is therefore a \textbf{similarity coefficient} in the sense of
Lesot et al.~{[}12{]} --- satisfying reflexivity and boundedness, with
conditional symmetry, but not constituting a metric distance. This
places it in the same class as Jaro-Winkler and other positionally
sensitive comparators used in record linkage.

\begin{center}\rule{0.5\linewidth}{0.5pt}\end{center}

\section{7. Experimental evaluation}\label{experimental-evaluation}

\subsection{7.1 Synthetic dataset}\label{synthetic-dataset}

To evaluate ILCSS against competing measures on the problem class for
which it is designed, we constructed a synthetic benchmark of 295 string
pairs covering the variation types described in Section 1. Starting from
22 base names drawn from five naming traditions --- Arabic (3- and
4-part chains), Hispanic (2-, 3-, and 4-part), Portuguese (2- and
3-part), and European 2-part (Italian, French, German) --- we generated:

\begin{itemize}
\tightlist
\item
  \textbf{All non-identity permutations} of each name's components (2
  permutations for 2-part, 5 for 3-part, 23 for 4-part);
\item
  \textbf{Component-drop variants}: removal of the first, middle, or
  last part from 3- and 4-part names;
\item
  \textbf{Insertion variants} for 2-part names: addition of a middle
  component or a trailing suffix;
\item
  \textbf{Negative (different-person) pairs}: each base name paired with
  three names from different entries in the set, spanning cross-cultural
  combinations.
\end{itemize}

The dataset contains 22 identity pairs, 207 same-person variant pairs,
and 66 different-person negative pairs (273 pairs excluding identities).
Ground truth is exact by construction. All names are in romanized Latin
script; Arabic-script orthographic variation is not tested here (see
Section 7.4).

\subsection{7.2 Algorithms compared}\label{algorithms-compared}

Four algorithms were evaluated on each pair:

{\def\LTcaptype{none} % do not increment counter
\begin{longtable}[]{@{}
  >{\raggedright\arraybackslash}p{(\linewidth - 2\tabcolsep) * \real{0.5000}}
  >{\raggedright\arraybackslash}p{(\linewidth - 2\tabcolsep) * \real{0.5000}}@{}}
\toprule\noalign{}
\begin{minipage}[b]{\linewidth}\raggedright
Algorithm
\end{minipage} & \begin{minipage}[b]{\linewidth}\raggedright
Description
\end{minipage} \\
\midrule\noalign{}
\endhead
\bottomrule\noalign{}
\endlastfoot
\textbf{ILCSS} & This work; threshold 3, positional penalty as described
in Section 3 \\
\textbf{NormLev} & 1 -- edit\_distance(s1, s2) / max( \\
\textbf{Jaro-Winkler} & Standard implementation (prefix scale p =
0.1) \\
\textbf{Ratcliff/Obershelp} & Recursive LCS; similarity = 2 $\times$ matched /
( \\
\end{longtable}
}

\subsection{7.3 Results}\label{results}

\textbf{Mean similarity scores by variation class} (identity pairs
excluded from aggregates):

{\def\LTcaptype{none} % do not increment counter
\begin{longtable}[]{@{}
  >{\raggedright\arraybackslash}p{(\linewidth - 10\tabcolsep) * \real{0.1667}}
  >{\raggedright\arraybackslash}p{(\linewidth - 10\tabcolsep) * \real{0.1667}}
  >{\raggedright\arraybackslash}p{(\linewidth - 10\tabcolsep) * \real{0.1667}}
  >{\raggedright\arraybackslash}p{(\linewidth - 10\tabcolsep) * \real{0.1667}}
  >{\raggedright\arraybackslash}p{(\linewidth - 10\tabcolsep) * \real{0.1667}}
  >{\raggedright\arraybackslash}p{(\linewidth - 10\tabcolsep) * \real{0.1667}}@{}}
\toprule\noalign{}
\begin{minipage}[b]{\linewidth}\raggedright
Variation class
\end{minipage} & \begin{minipage}[b]{\linewidth}\raggedright
ILCSS
\end{minipage} & \begin{minipage}[b]{\linewidth}\raggedright
NormLev
\end{minipage} & \begin{minipage}[b]{\linewidth}\raggedright
Jaro-Winkler
\end{minipage} & \begin{minipage}[b]{\linewidth}\raggedright
Ratcliff/Ob
\end{minipage} & \begin{minipage}[b]{\linewidth}\raggedright
N
\end{minipage} \\
\midrule\noalign{}
\endhead
\bottomrule\noalign{}
\endlastfoot
Identity & 1.000 & 1.000 & 1.000 & 1.000 & 22 \\
2-part swap & 0.757 & 0.110 & 0.574 & 0.517 & 9 \\
3-part permutations (all) & 0.788 & 0.402 & 0.789 & 0.640 & 45 \\
4-part permutations (all) & 0.781 & 0.433 & 0.813 & 0.639 & 92 \\
Drop one component & 0.654 & 0.683 & 0.877 & 0.809 & 43 \\
Insert / append & 0.823 & 0.796 & 0.949 & 0.886 & 18 \\
\textbf{Same person (all variants)} & \textbf{0.759} & \textbf{0.496} &
\textbf{0.823} & \textbf{0.691} & \textbf{207} \\
\textbf{Different person (negative)} & \textbf{0.030} & \textbf{0.231} &
\textbf{0.561} & \textbf{0.330} & \textbf{66} \\
\textbf{Gap $\Delta$ (same -- different)} & \textbf{0.729} & \textbf{0.265} &
\textbf{0.262} & \textbf{0.361} & --- \\
\end{longtable}
}

The most diagnostic figure is the gap between the mean scores for
same-person pairs and different-person pairs. ILCSS achieves $\Delta$ = 0.729,
more than twice the gap of Jaro-Winkler ($\Delta$ = 0.262) and Normalized
Levenshtein ($\Delta$ = 0.265), and substantially better than
Ratcliff/Obershelp ($\Delta$ = 0.361). The large gap reflects two complementary
behaviors: ILCSS scores reordered same-person pairs high ($\geq$ 0.75 across
all permutation types) while scoring unrelated name pairs close to zero
(mean 0.030), because the trigram threshold suppresses the short chance
matches between structurally unrelated names.

Jaro-Winkler achieves the highest mean score on same-person variants
(0.823) but also the highest mean on different-person pairs (0.561),
reducing the discriminative margin to the point where a fixed decision
threshold cannot separate the two classes reliably. Normalized
Levenshtein degrades sharply on permutations (0.110 for 2-part swaps)
because even identical character sets reordered within a string entail a
near-total character-level rewrite.

\textbf{Optimal decision threshold and F1 performance:}

{\def\LTcaptype{none} % do not increment counter
\begin{longtable}[]{@{}llllll@{}}
\toprule\noalign{}
Algorithm & Optimal $\theta$ & Precision & Recall & F1 & Accuracy \\
\midrule\noalign{}
\endhead
\bottomrule\noalign{}
\endlastfoot
\textbf{ILCSS} & \textbf{0.40} & \textbf{0.995} & \textbf{1.000} &
\textbf{0.998} & \textbf{0.996} \\
Ratcliff/Obershelp & 0.45 & 0.948 & 0.971 & 0.959 & 0.938 \\
Jaro-Winkler & 0.65 & 0.960 & 0.928 & 0.944 & 0.916 \\
NormLev & 0.25 & 0.884 & 0.879 & 0.881 & 0.821 \\
\end{longtable}
}

At threshold $\theta$ = 0.40, ILCSS produces a single false negative (one
same-person pair scored just below the threshold) and zero false
positives, for F1 = 0.998 and accuracy = 0.996. The precision curve for
ILCSS reaches 1.000 at $\theta$ = 0.40 and remains at 1.000 for all higher
thresholds, because no different-person pair scores above 0.40 in this
dataset. The recall curve remains above 0.96 down to $\theta$ = 0.55,
indicating robustness to threshold choice over a wide range.

Jaro-Winkler requires a substantially higher threshold (0.65) to achieve
acceptable precision, because many different-person pairs score between
0.50 and 0.65; at the same threshold, recall drops to 0.928, meaning 7\%
of same-person pairs are missed. Normalized Levenshtein achieves its
best F1 at the low threshold of 0.25, reflecting the algorithm's
inability to score permutation pairs above 0.50 --- a correct
classification regime but with poor separation and correspondingly lower
accuracy.

\subsection{7.4 Discussion}\label{discussion}

\textbf{Component drop.} For pairs where one name is a proper subset of
another (e.g., \texttt{"juan\ garcia\ lopez"} /
\texttt{"juan\ garcia"}), ILCSS gives a middle-range score (mean 0.664)
rather than a high one. This is the intended behavior: normalization by
the longer string penalizes missing components, preventing spurious
high-confidence matches when a name has been truncated. In a real
record-linkage system, such pairs would fall in a decision region
requiring additional evidence, not automatic acceptance.

\textbf{Positional penalty in context.} The penalty attenuates but does
not eliminate the positional signal. A fully reversed 3-part name scores
approximately 0.69--0.85, depending on component lengths; a fully
reversed 2-part name scores approximately 0.75. These scores are well
above the optimal threshold of 0.40, confirming that the penalty does
not harm recall for the reordering case.

\textbf{Jaro-Winkler on multi-part names.} Jaro-Winkler performs
surprisingly well on 3- and 4-part permutations (mean 0.771--0.820), but
the high scores on unrelated names (mean 0.561) make discrimination
difficult. The issue is structural: Jaro-Winkler's character-matching
window scales with string length, so for long multi-part names it finds
many character-level coincidences even between unrelated names sharing
common suffixes (e.g., \texttt{-ez}, \texttt{-ia}, \texttt{-er}). ILCSS
avoids this because the trigram threshold requires at least three
consecutive matching characters, which is a much stronger condition than
isolated character matches within a window.

\textbf{Limitations of the synthetic evaluation.} The dataset pairs are
generated by exact structural transformations; real-world name pairs
additionally exhibit typographic errors, transliteration variants,
partial abbreviations (initials), and combinations of multiple variation
types. The component-drop and insert tests cover only single-component
changes; chains of two or more simultaneous variations are not
represented. Furthermore, all names in the dataset are in romanized
Latin script; Arabic-script orthographic variation --- as catalogued for
example by Hassanat \& Altarawneh {[}26{]} --- is a distinct and
complementary challenge not addressed here. An evaluation on a real
annotated corpus of name pairs (see Section 8) would be required to
confirm these findings in operational conditions.

\begin{center}\rule{0.5\linewidth}{0.5pt}\end{center}

\section{8. Evaluation on real-world
data}\label{evaluation-on-real-world-data}

To complement the controlled synthetic evaluation (Section 7), we tested
ILCSS on two publicly available record-linkage benchmarks that contain
real name variation, albeit of a different character than the multi-part
reordering problem.

\subsection{8.1 DBLP-ACM author names}\label{dblp-acm-author-names}

The DBLP-ACM dataset (Leipzig Benchmark, CC BY 4.0 {[}27{]}) contains
2,224 matched pairs of publication records drawn from the DBLP and ACM
digital libraries, with a manually verified perfect mapping as ground
truth. For each matched pair we extracted the author name strings,
normalised to lower-case with initial letters removed, and compared the
DBLP author string against the ACM author string as a single character
sequence; 5 pairs were dropped because one record had an empty author
field, leaving 2,219 usable matched pairs. Non-matched pairs (N = 222)
were sampled by misaligning indices across the matched set.

The dataset is not designed for within-name component reordering: it
contains full author lists (multiple names separated by commas), and
variation arises from abbreviation styles, typographic errors, and
inconsistent author-list ordering across the two libraries. Nonetheless,
the matched pairs exhibit the same structural challenge ILCSS addresses:
long strings with identical content in different positional
arrangements.

{\def\LTcaptype{none} % do not increment counter
\begin{longtable}[]{@{}
  >{\raggedright\arraybackslash}p{(\linewidth - 12\tabcolsep) * \real{0.1429}}
  >{\raggedright\arraybackslash}p{(\linewidth - 12\tabcolsep) * \real{0.1429}}
  >{\raggedright\arraybackslash}p{(\linewidth - 12\tabcolsep) * \real{0.1429}}
  >{\raggedright\arraybackslash}p{(\linewidth - 12\tabcolsep) * \real{0.1429}}
  >{\raggedright\arraybackslash}p{(\linewidth - 12\tabcolsep) * \real{0.1429}}
  >{\raggedright\arraybackslash}p{(\linewidth - 12\tabcolsep) * \real{0.1429}}
  >{\raggedright\arraybackslash}p{(\linewidth - 12\tabcolsep) * \real{0.1429}}@{}}
\toprule\noalign{}
\begin{minipage}[b]{\linewidth}\raggedright
Measure
\end{minipage} & \begin{minipage}[b]{\linewidth}\raggedright
Same person (N=2219)
\end{minipage} & \begin{minipage}[b]{\linewidth}\raggedright
Different (N=222)
\end{minipage} & \begin{minipage}[b]{\linewidth}\raggedright
Gap
\end{minipage} & \begin{minipage}[b]{\linewidth}\raggedright
Best thr.
\end{minipage} & \begin{minipage}[b]{\linewidth}\raggedright
F1
\end{minipage} & \begin{minipage}[b]{\linewidth}\raggedright
Accuracy
\end{minipage} \\
\midrule\noalign{}
\endhead
\bottomrule\noalign{}
\endlastfoot
\textbf{ILCSS} & \textbf{0.783} & \textbf{0.025} & \textbf{0.758} &
\textbf{0.15} & \textbf{0.994} & \textbf{0.990} \\
Ratcliff/Ob & 0.716 & 0.223 & 0.493 & 0.30 & 0.984 & 0.971 \\
Jaro-Winkler & 0.854 & 0.548 & 0.306 & 0.60 & 0.971 & 0.948 \\
NormLev & 0.569 & 0.185 & 0.384 & 0.05 & 0.950 & 0.904 \\
\end{longtable}
}

ILCSS achieves the largest gap between matched and non-matched pairs
(0.758 vs.~0.306--0.493 for the other algorithms), and reaches F1 =
0.994 at threshold 0.15 --- lower than the synthetic-data optimum of
0.40, reflecting that some matched author strings differ substantially
in length due to abbreviation. Jaro-Winkler's gap is the smallest
(0.306) because its character-window mechanism matches many characters
between unrelated author lists that share common name fragments.

We identified a small subset of 1 pair (out of 2,219) where the matched
records contain the exact same author-name tokens in a different order
--- confirming that this variation type exists in the real data, though
at low frequency in this particular benchmark. For that pair: ILCSS =
0.780, NormLev = 0.548, Jaro-Winkler = 0.865, Ratcliff/Ob = 0.715.

\subsection{8.2 FEBRL given-name / surname
swap}\label{febrl-given-name-surname-swap}

The FEBRL dataset1 {[}28{]} is a synthetic record-linkage benchmark with
500 matched pairs generated with controlled noise: field corruption,
typographic errors, and --- relevant to ILCSS --- a given-name/surname
field swap applied to approximately 7\% of matched pairs. We
concatenated the \texttt{given\_name} and \texttt{surname} fields into a
single name string before comparison.

{\def\LTcaptype{none} % do not increment counter
\begin{longtable}[]{@{}
  >{\raggedright\arraybackslash}p{(\linewidth - 12\tabcolsep) * \real{0.1429}}
  >{\raggedright\arraybackslash}p{(\linewidth - 12\tabcolsep) * \real{0.1429}}
  >{\raggedright\arraybackslash}p{(\linewidth - 12\tabcolsep) * \real{0.1429}}
  >{\raggedright\arraybackslash}p{(\linewidth - 12\tabcolsep) * \real{0.1429}}
  >{\raggedright\arraybackslash}p{(\linewidth - 12\tabcolsep) * \real{0.1429}}
  >{\raggedright\arraybackslash}p{(\linewidth - 12\tabcolsep) * \real{0.1429}}
  >{\raggedright\arraybackslash}p{(\linewidth - 12\tabcolsep) * \real{0.1429}}@{}}
\toprule\noalign{}
\begin{minipage}[b]{\linewidth}\raggedright
Measure
\end{minipage} & \begin{minipage}[b]{\linewidth}\raggedright
Same person (N=500)
\end{minipage} & \begin{minipage}[b]{\linewidth}\raggedright
Different (N=100)
\end{minipage} & \begin{minipage}[b]{\linewidth}\raggedright
Gap
\end{minipage} & \begin{minipage}[b]{\linewidth}\raggedright
Best thr.
\end{minipage} & \begin{minipage}[b]{\linewidth}\raggedright
F1
\end{minipage} & \begin{minipage}[b]{\linewidth}\raggedright
Accuracy
\end{minipage} \\
\midrule\noalign{}
\endhead
\bottomrule\noalign{}
\endlastfoot
\textbf{ILCSS} & \textbf{0.825} & \textbf{0.010} & \textbf{0.816} &
\textbf{0.15} & \textbf{0.989} & \textbf{0.982} \\
Ratcliff/Ob & 0.874 & 0.268 & 0.606 & 0.40 & 0.987 & 0.978 \\
Jaro-Winkler & 0.905 & 0.512 & 0.393 & 0.65 & 0.945 & 0.913 \\
NormLev & 0.811 & 0.187 & 0.625 & 0.30 & 0.944 & 0.910 \\
\end{longtable}
}

For the 35 swapped pairs specifically (same tokens, first/last position
exchanged), ILCSS scores mean 0.758, maintaining good recall at
threshold 0.15. NormLev degrades sharply to 0.171 on these pairs ---
consistent with the synthetic results --- while Jaro-Winkler (0.621) and
Ratcliff/Ob (0.531) both score below the 0.65 and 0.40 thresholds
respectively, meaning they would misclassify a fraction of the swapped
pairs as non-matches.

\subsection{8.3 Discussion}\label{discussion-1}

Across both real datasets, ILCSS consistently achieves the largest gap
between matched and non-matched name pairs. The F1 optima on real data
(0.15 for DBLP-ACM and FEBRL) are lower than on the synthetic data
(0.40), because real matched records contain noise that reduces scores
for genuine matches. This suggests that a threshold of approximately
0.40 is appropriate when names are relatively clean (administrative
databases with standardized input), while a threshold of 0.15--0.25 is
more appropriate when names are subject to abbreviation, partial
records, or transcription errors.

We note a gap in the available public benchmarks: no dataset with ground
truth currently covers the full problem stated in Section 1 --- namely,
multi-part personal names from culturally diverse sources where the same
individual's components appear in systematically different orders. The
LREC 2008 dataset (Arehart \& Miller) {[}29{]} was specifically designed
for this problem and includes Arabic names with segmentation variation,
but has not been publicly released. Filling this benchmark gap is an
important open task for the record-linkage community.

\begin{center}\rule{0.5\linewidth}{0.5pt}\end{center}

\section{References}\label{references}

{[}1{]} C. N. Alberga, ``String Similarity and Misspellings,''
\emph{Communications of the ACM}, Vol. 10, No.~5, May 1967,
pp.~302--313. DOI: 10.1145/363282.363371. \emph{(Cites H. Baskin and L.
D. Selfridge, IBM Research Center, as private communications {[}refs. 1
and 3 therein{]} --- the iterative longest-common-substring method was
never independently published by either author.)}

{[}2{]} J. W. Ratcliff and D. E. Metzener, ``Pattern Matching: The
Gestalt Approach,'' \emph{Dr.~Dobb's Journal}, Vol. 13, No.~7, July
1988, p.~46.

{[}3{]} R. A. Wagner and M. J. Fischer, ``The String-to-String
Correction Problem,'' \emph{Journal of the ACM}, Vol. 21, No.~1, January
1974, pp.~168--173. DOI: 10.1145/321796.321811.

{[}4{]} D. Gusfield, \emph{Algorithms on Strings, Trees and Sequences:
Computer Science and Computational Biology}, Cambridge University Press,
1997. ISBN: 978-0-521-58519-4.

{[}5{]} D. S. Hirschberg, ``A Linear Space Algorithm for Computing
Maximal Common Subsequences,'' \emph{Communications of the ACM}, Vol.
18, No.~6, June 1975, pp.~341--343. DOI: 10.1145/360825.360861.

{[}6{]} G. Navarro, ``A Guided Tour to Approximate String Matching,''
\emph{ACM Computing Surveys}, Vol. 33, No.~1, March 2001, pp.~31--88.
DOI: 10.1145/375360.375365.

{[}7{]} R. Cox, ``Regular Expression Matching with a Trigram Index,''
2012. https://swtch.com/\textasciitilde rsc/regexp/regexp4.html

{[}8{]} E. Ukkonen, ``Approximate String-Matching with q-Grams and
Maximal Matches,'' \emph{Theoretical Computer Science}, Vol. 92, No.~1,
1992, pp.~191--211. DOI: 10.1016/0304-3975(92)90143-4.

{[}9{]} P. Willett, ``Document Retrieval Experiments Using Indexing
Vocabularies of Varying Size. II. Hashing, Truncation, Digram and
Trigram Encoding of Index Terms,'' \emph{Journal of Documentation}, Vol.
35, No.~4, 1979, pp.~296--305. DOI: 10.1108/eb026684.

{[}10{]} R. C. Angell, G. E. Freund, P. Willett, ``Automatic Spelling
Correction Using a Trigram Similarity Measure,'' \emph{Information
Processing \& Management}, Vol. 19, No.~4, 1983, pp.~255--261. DOI:
10.1016/0306-4573(83)90022-5.

{[}11{]} S. Santini and R. Jain, ``Similarity Measures,'' \emph{IEEE
Transactions on Pattern Analysis and Machine Intelligence}, Vol. 21,
No.~9, September 1999, pp.~871--883. DOI: 10.1109/34.790428.

{[}12{]} M.-J. Lesot, M. Rifqi, H. Benhadda, ``Similarity Measures for
Binary and Numerical Data: A Survey,'' \emph{International Journal of
Knowledge Engineering and Soft Data Paradigms}, Vol. 1, No.~1, 2009,
pp.~63--84. DOI: 10.1504/IJKESDP.2009.021985.

{[}13{]} G. Kondrak, ``N-Gram Similarity and Distance,'' in \emph{String
Processing and Information Retrieval (SPIRE 2005)}, LNCS vol.~3772,
Springer, 2005, pp.~115--126. DOI: 10.1007/11575832\_13.

{[}14{]} Li Yujian and Liu Bo, ``A Normalized Levenshtein Distance
Metric,'' \emph{IEEE Transactions on Pattern Analysis and Machine
Intelligence}, Vol. 29, No.~6, June 2007, pp.~1091--1095. DOI:
10.1109/TPAMI.2007.1078.

{[}15{]} M. A. Jaro, ``Advances in Record-Linkage Methodology as Applied
to Matching the 1985 Census of Tampa, Florida,'' \emph{Journal of the
American Statistical Association}, Vol. 84, No.~406, June 1989,
pp.~414--420. DOI: 10.1080/01621459.1989.10478785.

{[}16{]} W. E. Winkler, ``String Comparator Metrics and Enhanced
Decision Rules in the Fellegi-Sunter Model of Record Linkage,''
\emph{Proceedings of the Section on Survey Research Methods}, American
Statistical Association, 1990, pp.~354--359. ERIC: ED325505.

{[}17{]} W. W. Cohen, P. Ravikumar, S. E. Fienberg, ``A Comparison of
String Distance Metrics for Name-Matching Tasks,'' \emph{Proceedings of
the IJCAI-2003 Workshop on Information Integration on the Web
(IIWeb-2003)}, 2003, pp.~73--78.

{[}18{]} I. P. Fellegi and A. B. Sunter, ``A Theory for Record
Linkage,'' \emph{Journal of the American Statistical Association}, Vol.
64, No.~328, 1969, pp.~1183--1210. DOI: 10.1080/01621459.1969.10501049.

{[}19{]} P. Christen, \emph{Data Matching: Concepts and Techniques for
Record Linkage, Entity Resolution, and Duplicate Detection}, Springer,
2012. ISBN: 978-3-642-31163-5. DOI: 10.1007/978-3-642-31164-2.

{[}20{]} P. Christen, ``A Comparison of Personal Name Matching:
Techniques and Practical Issues,'' \emph{Proceedings of the Sixth IEEE
International Conference on Data Mining Workshops (ICDMW 2006)}, 2006,
pp.~290--294. DOI: 10.1109/ICDMW.2006.2.

{[}21{]} A. J. Lait and B. Randell, ``An Assessment of Name Matching
Algorithms,'' Technical Report No.~550, Department of Computing Science,
University of Newcastle upon Tyne, 1996.

{[}22{]} R. Ruiz-P\'erez, E. Delgado L\'opez-C\'ozar, E. Jim\'enez-Contreras,
``Spanish Personal Name Variations in National and International
Biomedical Databases: Implications for Information Retrieval and
Bibliometric Studies,'' \emph{Journal of the Medical Library
Association}, Vol. 90, No.~4, 2002, pp.~411--430. PMCID: PMC128958.

{[}23{]} D. R. Wilson, ``Name Standardization for Genealogical Record
Linkage,'' \emph{Proceedings of the Family History Technology Workshop
(FHTW 2005)}, Brigham Young University, 2005.

{[}24{]} T. El-Shishtawy, ``Linking Databases using Matched Arabic
Names,'' \emph{International Journal of Computational Linguistics \&
Chinese Language Processing}, Vol. 19, No.~1, 2014, pp.~33--54. ACL:
O14-2003.

{[}25{]} A. Freeman, S. Condon, C. Ackerman, ``Cross Linguistic Name
Matching in English and Arabic,'' \emph{Proceedings of the Human
Language Technology Conference of the NAACL, Main Conference}, 2006,
pp.~471--478. ACL: N06-1060.

{[}26{]} A. B. A. Hassanat and G. A. Altarawneh, ``Rule-and
Dictionary-based Solution for Variations in Written Arabic Names in
Social Networks, Big Data, Accounting Systems and Large Databases,''
\emph{Research Journal of Applied Sciences, Engineering and Technology},
Vol. 8, No.~14, 2014, pp.~1630--1638. arXiv: 1502.05441.

{[}27{]} E. K\"opcke, A. Thor, E. Rahm, ``Evaluation of Entity Resolution
Approaches on Real-World Match Problems,'' \emph{Proceedings of the VLDB
Endowment}, Vol. 3, No.~1, 2010, pp.~484--493. Dataset: DBLP-ACM,
available at
https://dbs.uni-leipzig.de/research/projects/object\_matching/benchmark\_datasets\_for\_entity\_resolution
(CC BY 4.0).

{[}28{]} P. Christen and K. Goiser, ``Quality and Complexity Measures
for Data Linkage and Deduplication,'' in F. Guillet and H. Hamilton
(eds.), \emph{Quality Measures in Data Mining}, Springer, 2007,
pp.~127--151. FEBRL dataset available via Python Record Linkage Toolkit
(https://github.com/J535D165/recordlinkage, BSD-3-Clause).

{[}29{]} M. Arehart and K. J. Miller, ``A Ground Truth Dataset for
Matching Culturally Diverse Romanized Person Names,'' \emph{Proceedings
of LREC 2008}, Marrakech, 2008, pp.~2064--2069.

\end{document}
